% !TEX root = ../main.tex
% !TEX program = lualatex
\documentclass[../main.tex]{subfiles}
\IfSubfilesClassLoaded{\externaldocument{\subfix{../build/main}}}{}
% =============================================================================
% CHAPTER 12A: SOURCE SELECTION PROCESS
% DESCRIPTION: Source selection governance, interactions, and negotiations.
% =============================================================================
\begin{document}
\IfSubfilesClassLoaded{\chapter{EDO Study Guide}}{}

% === SOURCE SELECTION PROCESS (EDO 3.2.3) ===
\section{Source Selection Process}

\subsection{Overview}
\begin{itemize}
  \item Formal source selection is used for high-dollar defense acquisitions to achieve best value in cost, quality, and timeliness~\autocite{edo-3-2-3-source-selection-2025}.
  \item Source selection participants may be sequestered and required to sign a non-disclosure agreement~\autocite{edo-3-2-3-source-selection-2025}.
  \item The process is guided by the \ac{ssplan} and the solicitation (\ac{rfp} or \ac{ifb})~\autocite{edo-3-2-3-source-selection-2025}.
\end{itemize}

\subsection{Source Selection Plan}
\begin{description}
  \item[\textbf{Purpose.}] The \ac{ssplan} details how the Government will solicit and evaluate proposals and execute source selection~\autocite{edo-3-2-3-source-selection-2025}.
  \item[\textbf{Origins and approval.}] It originates from the acquisition strategy, is prepared by the \ac{ko} with the \ac{ipt}, and is approved by the \ac{ssa}~\autocite{edo-3-2-3-source-selection-2025}.
  \item[\textbf{Contents.}] Description of the procurement, source selection charter, team organization, best-value approach, evaluation factors and rating scheme, and milestone schedule~\autocite{edo-3-2-3-source-selection-2025}.
\end{description}

\subsection{Basis for Award and Best Value}
\begin{itemize}
  \item Sealed bidding focuses on responsiveness and responsibility~\autocite{edo-3-2-3-source-selection-2025}.
  \item Negotiated acquisitions select best value, defined as the outcome that provides the greatest overall benefit in response to the requirement~\autocite{edo-3-2-3-source-selection-2025}.
  \item Best-value policies include \ac{lpta} and tradeoff processes~\autocite{edo-3-2-3-source-selection-2025}.
\end{itemize}

\subsection{Lowest Price Technically Acceptable}
\begin{itemize}
  \item Best value results from the technically acceptable proposal with the lowest evaluated cost or price~\autocite{edo-3-2-3-source-selection-2025}.
  \item Once technical acceptability is established, proposals are ranked on cost or price only~\autocite{edo-3-2-3-source-selection-2025}.
  \item Tradeoffs are not permitted~\autocite{edo-3-2-3-source-selection-2025}.
\end{itemize}

\subsection{Tradeoff Process}
\begin{itemize}
  \item Best value is a balance of technical merit, performance risk, management capability, and price factors~\autocite{edo-3-2-3-source-selection-2025}.
  \item Used when requirements are less defined, significant development work is needed, or performance risk is high~\autocite{edo-3-2-3-source-selection-2025}.
  \item The solicitation must state evaluation factors and their relative importance, and the rationale for selection must be documented~\autocite{edo-3-2-3-source-selection-2025}.
\end{itemize}

\subsection{Source Selection Steps}
\begin{description}
  \item[\textbf{Common start.}] Compare each proposal to the \ac{rfp} requirements~\autocite{edo-3-2-3-source-selection-2025}.
  \item[\textbf{Tradeoff path.}] Perform comparative analysis, determine whether prices are fair and reasonable, and select the source for award~\autocite{edo-3-2-3-source-selection-2025}.
  \item[\textbf{\ac{lpta} path.}] Identify technically acceptable sources, then select the lowest-priced acceptable offer for award~\autocite{edo-3-2-3-source-selection-2025}.
\end{description}

\subsection{Source Selection Organization}
\begin{itemize}
  \item The \ac{ssa} is the decision authority; the \ac{ssac} and advisors compare proposals and provide recommendations~\autocite{edo-3-2-3-source-selection-2025}.
  \item The \ac{sseb} evaluates proposals and is typically organized into technical, management, past performance, cost/pricing, and small business teams as needed~\autocite{edo-3-2-3-source-selection-2025}.
\end{itemize}

\subsection{Roles and Responsibilities}
\begin{description}
  \item[\textbf{\ac{pco}/\ac{ko}.}] Ensures compliance with procurement laws, reviews proposal completeness, manages communications, and awards the contract~\autocite{edo-3-2-3-source-selection-2025}.
  \item[\textbf{\ac{ssa}.}] Oversees integrity of the process, appoints qualified personnel, and selects the best-value proposal~\autocite{edo-3-2-3-source-selection-2025}.
  \item[\textbf{\ac{ssac}.}] Reviews proposal analyses after \ac{sseb} evaluation and recommends the best-value source; functions may be assigned to the \ac{sseb} when a separate council is not used~\autocite{edo-3-2-3-source-selection-2025}.
  \item[\textbf{\ac{sseb}.}] Evaluates proposals against Section M criteria and rating standards, identifies strengths and weaknesses, documents deficiencies, and drafts discussion questions~\autocite{edo-3-2-3-source-selection-2025}.
\end{description}

\subsection{Interactions with Contractors}
\begin{description}
  \item[\textbf{Deficiencies.}] Material failures to meet minimum requirements or combinations of significant weaknesses that increase performance risk; must be corrected, and correction requires proposal revision~\autocite{edo-3-2-3-source-selection-2025}.
  \item[\textbf{Competitive range.}] The most highly rated offerors selected for further discussions~\autocite{edo-3-2-3-source-selection-2025}.
  \item[\textbf{Clarifications.}] Used when awarding without discussions; limited to minor or clerical errors and adverse past performance issues~\autocite{edo-3-2-3-source-selection-2025}.
  \item[\textbf{Communications.}] Used when inclusion in the competitive range is uncertain, to enhance Government understanding without proposal revision~\autocite{edo-3-2-3-source-selection-2025}.
  \item[\textbf{Discussions.}] Conducted after establishing the competitive range; offerors may revise proposals and negotiations may occur~\autocite{edo-3-2-3-source-selection-2025}.
\end{description}

\subsection{Guidance on Exchanges}
\begin{itemize}
  \item Do not favor one offeror over another or reveal another offeror's technical solution or price~\autocite{edo-3-2-3-source-selection-2025}.
  \item Conduct meaningful discussions with all offerors in the competitive range, addressing all deficiencies and significant weaknesses~\autocite{edo-3-2-3-source-selection-2025}.
  \item Allow reasonable time for proposal revision and document all discussions (oral or written)~\autocite{edo-3-2-3-source-selection-2025}.
\end{itemize}

\subsection{Fact-Finding and Analysis}
\begin{itemize}
  \item Fact-finding prepares the Government's pre-negotiation position and relies on Government research and controlled interactions that do not permit proposal revision~\autocite{edo-3-2-3-source-selection-2025}.
  \item Typical analyses include technical analysis, cost analysis, price analysis, and cost-realism analysis (required for cost-reimbursement contracts)~\autocite{edo-3-2-3-source-selection-2025}.
  \item External support can include \ac{dcma}, \ac{dcaa}, and field activities~\autocite{edo-3-2-3-source-selection-2025}.
\end{itemize}

\subsection{Preparing for Fact-Finding}
\begin{enumerate}
  \item Review the contracting situation and relevant legal and contractual terms~\autocite{edo-3-2-3-source-selection-2025}.
  \item Identify team members who will support analysis and negotiations~\autocite{edo-3-2-3-source-selection-2025}.
  \item List open issues, conduct research through evaluation teams, and develop the pre-negotiation position using cost and price techniques~\autocite{edo-3-2-3-source-selection-2025}.
  \item Draft questions and submit them in writing to the \ac{ko}, who forwards them to offerors~\autocite{edo-3-2-3-source-selection-2025}.
\end{enumerate}

\subsection{Negotiation Preparation and Objectives}
\begin{itemize}
  \item Be familiar with the solicitation and proposals, submit questions in advance, and align the team on acceptable trade-offs~\autocite{edo-3-2-3-source-selection-2025}.
  \item Handle administrative details early (agenda, assignments, facilities)~\autocite{edo-3-2-3-source-selection-2025}.
  \item The \ac{ko} establishes pre-negotiation objectives using field pricing reports, \ac{dcaa} audits, technical analysis, \ac{igce} data, and price histories, with management review per policy~\autocite{edo-3-2-3-source-selection-2025}.
\end{itemize}

\subsection{Negotiation Skills}
\begin{description}
  \item[\textbf{Do.}] Identify and prioritize discussion items, verify contractor authority, listen and probe for detail, draw estimate bases into the open, document action items, and keep thorough notes~\autocite{edo-3-2-3-source-selection-2025}.
  \item[\textbf{Do not.}] Reveal Government findings or specific numbers, coach the contractor, negotiate how to perform tasks, or let issues go unresolved~\autocite{edo-3-2-3-source-selection-2025}.
\end{description}

\subsection{Keep the End Goal in Mind}
\begin{itemize}
  \item The goal of negotiations is an agreement that enhances the quality of the product or service delivered to the Government~\autocite{edo-3-2-3-source-selection-2025}.
  \item Ensure negotiation activities support that end state~\autocite{edo-3-2-3-source-selection-2025}.
\end{itemize}

\subsection{Contract Award and Protests}
\begin{itemize}
  \item The \ac{ssa} signs the source selection decision document; the \ac{ko} incorporates negotiated changes, obtains legal review, and signs the contract~\autocite{edo-3-2-3-source-selection-2025}.
  \item Only the \ac{ko} can award; the contract is binding only after signature and distribution, and the \ac{ko} must ensure price reasonableness and protection of Government interests~\autocite{edo-3-2-3-source-selection-2025}.
  \item Unsuccessful offerors are notified and may request a debriefing on strengths and weaknesses~\autocite{edo-3-2-3-source-selection-2025}.
  \item Protests are written objections to award, cancellation, or termination; agency protests are governed by Executive Order 12979, with Alternative Dispute Resolution preferred and escalation to the \ac{cofc} if needed~\autocite{edo-3-2-3-source-selection-2025}.
  \item Protests generally suspend contract performance unless an urgent need exists~\autocite{edo-3-2-3-source-selection-2025}.
\end{itemize}

\subsection{Summary}
\begin{itemize}
  \item Guiding documents: \ac{ssplan} and the \ac{rfp} or \ac{ifb}~\autocite{edo-3-2-3-source-selection-2025}.
  \item Formal steps: compare proposals to the \ac{rfp}, perform comparative analysis, determine fair and reasonable price, and select the source~\autocite{edo-3-2-3-source-selection-2025}.
  \item \ac{sseb} responsibilities: evaluate against Section M factors, document strengths and weaknesses, and draft discussion questions~\autocite{edo-3-2-3-source-selection-2025}.
  \item Fact-finding questions are communicated in writing through the \ac{ko}~\autocite{edo-3-2-3-source-selection-2025}.
  \item Successful negotiation steps: prepare, conduct fact-finding and analysis, set pre-negotiation objectives, ensure negotiating skills, and keep the end goal in mind~\autocite{edo-3-2-3-source-selection-2025}.
\end{itemize}

%====================
% End of file
%====================
\IfSubfilesClassLoaded{
  \RenewDocumentCommand{\entryname}{}{\textbf{\color{Modern} Acronym}}
  \RenewDocumentCommand{\descriptionname}{}{\textbf{\color{Modern} Definition}}
  \printnoidxglossary[
    type=\acronymtype,
    title=Acronyms,
    style=long-booktabs
  ]}{}
\end{document}
